\anexo{Prompts Utilizados na Estratégia Map-Reduce com Divisão de Tarefas}
\label{an:prompts-multi-map-reduce}

Este anexo apresenta os \textit{prompts} utilizados na abordagem baseada no 
paradigma \textit{Map-Reduce} com divisão de tarefas, adotada neste trabalho 
para a extração estruturada de informações a partir das transcrições das reuniões.

Nessa estratégia, a fase de \textbf{Map} é composta por múltiplos agentes 
especializados — um para cada tipo de informação relevante:  
(i) Participantes,  
(ii) Pautas,  
(iii) Temas discutidos,  
(iv) Deliberações.

Cada agente possui seu próprio \textit{prompt}, processando de forma independente 
cada segmento da transcrição.  
Na fase de \textbf{Reduce}, os resultados parciais de cada agente são consolidados 
separadamente, produzindo sínteses limpas para cada categoria antes da montagem 
final da ata.

% ============================================================ %
\section*{Prompts de Mapeamento (Map)}
\addcontentsline{toc}{subsection}{Prompts de Mapeamento (Map)}

A seguir, apresentam-se os \textit{prompts} de cada tarefa especializada utilizada 
na fase Map. Cada um deles é aplicado de forma independente a todos os segmentos 
de transcrição.

% ----------------------------- %
\subsection*{Task: FindParticipants}
\begin{lstlisting}[language=Python]
Você é um agente de IA responsável por identificar os participantes de uma 
reunião que ocorreu no Tribunal Regional Eleitoral do Ceará.

Sua única tarefa é identificar o nome completo e o cargo/função (se possível) 
de todos os participantes, convidados e ausentes mencionados no segmento 
da transcrição apresentado.

MANDATÓRIO: O RESULTADO DEVE SER INTEIRAMENTE EM PORTUGUÊS BRASILEIRO.

A saída deve ser uma lista simples no formato:
[Nome Completo] - [Cargo/Função]

--- INÍCIO DA TRANSCRIÇÃO ---
{chunk_text}
--- FIM DA TRANSCRIÇÃO ---
\end{lstlisting}

% ----------------------------- %
\subsection*{Task: FindPautas}
\begin{lstlisting}[language=Python]
Você é um agente de IA responsável por identificar as pautas discutidas 
durante uma reunião realizada no Tribunal Regional Eleitoral do Ceará.

Sua tarefa é listar cada pauta mencionada no segmento e, quando possível, 
associá-la ao participante ou órgão que a propôs.

MANDATÓRIO: O RESULTADO DEVE SER INTEIRAMENTE EM PORTUGUÊS BRASILEIRO.

A saída deve estar no formato:
[Título] - [Descrição da pauta]

--- INÍCIO DA TRANSCRIÇÃO ---
{chunk_text}
--- FIM DA TRANSCRIÇÃO ---
\end{lstlisting}

% ----------------------------- %
\subsection*{Task: FindDiscucoes}
\begin{lstlisting}[language=Python]
Você é um agente de IA responsável por identificar os temas discutidos em um 
segmento de transcrição de reunião do Tribunal Regional Eleitoral do Ceará.

Sua tarefa é identificar os temas e apresentar uma breve descrição que represente
o contexto e os pontos debatidos.

MANDATÓRIO: O RESULTADO DEVE SER INTEIRAMENTE EM PORTUGUÊS BRASILEIRO.

Formato da saída:
[Tema identificado] - [Descrição detalhada]

--- INÍCIO DA TRANSCRIÇÃO ---
{chunk_text}
--- FIM DA TRANSCRIÇÃO ---
\end{lstlisting}

% ----------------------------- %
\subsection*{Task: FindDeliberacoes}
\begin{lstlisting}[language=Python]
Você é um agente de IA responsável por identificar as deliberações tomadas em 
uma reunião do Tribunal Regional Eleitoral do Ceará.

Sua tarefa é identificar as ações deliberadas, o tema relacionado e, sempre que 
possível, indicar o responsável pela execução da ação.

MANDATÓRIO: O RESULTADO DEVE SER INTEIRAMENTE EM PORTUGUÊS BRASILEIRO.

--- INÍCIO DA TRANSCRIÇÃO ---
{chunk_text}
--- FIM DA TRANSCRIÇÃO ---
\end{lstlisting}

% ============================================================ %
\section*{Prompt de Redução (Reduce)}
\addcontentsline{toc}{subsection}{Prompt de Redução (Reduce)}

A fase Reduce é executada separadamente para cada agente especializado.  
Cada tarefa recebe todos os resultados brutos produzidos na fase Map 
e consolida essas informações em uma síntese final única.

\begin{lstlisting}[language=Python]
Você é um agente de consolidação de dados. Sua tarefa é unificar e limpar 
as informações extraídas de múltiplos segmentos de transcrição para a 
categoria '{task_display_name}'.

Instruções:
1. Remova duplicatas e informações conflitantes.
2. Consolide os itens em um único texto coeso.
3. Preserve o estilo e o conteúdo original.
4. NÃO adicione informações externas.
5. O texto final deve estar inteiramente em Português brasileiro.

DADOS BRUTOS:
{dados_brutos}
\end{lstlisting}